\section{Conclusion}
\label{sec:conclusion}
This report demonstrated the design, implementation, and evaluation of a modular ZK-Rollup benchmarking prototype, aimed at addressing the critical scalability challenges inherent in Ethereum Layer-1. Unlike production-grade rollups—which prioritize EVM equivalence and reliability at the cost of experimental flexibility—this prototype was explicitly architected to isolate and evaluate distinct performance bottlenecks under controlled, synthetic conditions.

By cleanly separating the sequencing, execution, proving, and submission stages into Rust microservices and providing a Python data-tools pipeline, the framework allowed for a systematic analysis of the fundamental throughput--latency--cost trade-offs in Layer-2 scaling. The empirical observations yielded several key conclusions:
\begin{enumerate}
    \item \textbf{The Batching Dilemma:} Larger batch sizes significantly reduce the amortized Layer-1 gas cost per transaction, but directly penalize early-arriving users by increasing queueing latency (delaying soft finality).
    \item \textbf{Data Availability Dynamics:} Transitioning from standard transaction \texttt{calldata} to EIP-4844 blobs drastically reduces the L1 execution gas footprint. This structural efficiency gain allows rollups to operate more cost-effectively, fundamentally shifting the profitability curve of Layer-2 submission.
    \item \textbf{Prover Fixed Costs:} Using real RISC Zero STARK proofs demonstrated a high fixed wall-clock setup cost (~75 seconds) and execution step-functions that heavily penalize small batches. Furthermore, through parameterized mock proving, the research demonstrated that if proof generation is significantly delayed, the system becomes strictly compute-bound, creating severe downstream backpressure.
\end{enumerate}

Ultimately, the prototype successfully serves its purpose as a reproducible research framework, validating that scalability optimizations cannot be achieved in isolation. Enhancing one dimension (e.g., cost via large batches) inevitably impacts others (e.g., latency or computational overhead).
